\documentclass[10pt,a4paper]{amsart}
\usepackage[margin=19mm]{geometry}
\usepackage{amsmath,amssymb,mathtools}
\usepackage[scr=rsfs,cal=euler]{mathalfa}
\usepackage{xfrac}
\usepackage[unicode,hidelinks,bookmarks=false]{hyperref}
\newcommand{\ie}{i.e.\ }
\newcommand{\cf}{cf.\ }
\newcommand{\resp}{respectively\ }
\newcommand{\Subsection}{\subsection}
\providecommand{\nobreakdash}{\nobreak\hbox{-}}
\setlength{\emergencystretch}{2em}
\multlinegap0pt
\usepackage{bm}
\usepackage{bbm}
\usepackage{dsfont}
\usepackage{xspace}
\usepackage{cancel}
\usepackage{mathtools}
\usepackage{amsthm}
\makeatletter
\@ifundefined{theorem}{\newtheorem{theorem}{Theorem}}{}
\makeatother
\makeatletter
\expandafter\ifx\csname dedication\endcsname\relax
\newenvironment{dedication}
    {\vspace{-10ex}\begin{quotation}\begin{center}\begin{em}}
    {\par\end{em}\end{center}\end{quotation}}
\fi
\makeatother
\title[Some spectral properties and convergence of the $ (A,q)$-num]{Some spectral properties and convergence of the $ (A,q)$-numerical radius and $ (A,q)$-Crawford number}
\author{Pembe Ipek Al, Zameddin I. Ismailov, Fuad Kittaneh, Satyajit Sahoo}
\date{Source: arXiv:2505.16924v1 (CC BY 4.0)}
\begin{document}
\maketitle

\noindent\emph{Source and licence.} Some spectral properties and convergence of the $ (A,q)$-numerical radius and $ (A,q)$-Crawford number. Version arXiv:2505.16924v1 (CC BY 4.0), distributed under the Creative Commons Attribution 4.0 International licence, as recorded in the arXiv licence metadata for this exact version and shown on the abstract page https://arxiv.org/abs/2505.16924v1. Full rights notice: licenses/arxiv-2505.16924-CC-BY-4.0.txt.\par\smallskip

\noindent\textit{Editorial scope.} Counted unit: the Theorem whose source label is \texttt{thm1} in the article (source0.tex lines 157--170, source label \texttt{thm1}), delivered together with the article's own complete original proof of it (source0.tex lines 171--240, one \texttt{proof} environment of 70 lines). No mathematics is written, repaired or re-derived: statement and proof are the article's own text line for line. Required context retained uncounted: none. Every definition, display and label of those lines is retained unchanged. Omitted from this package and not redistributed: 1--156, 241--853. Nothing inside a retained excerpt is omitted or altered: every retained body is delivered exactly as the article's lines are written, so no omission receipt of any kind accompanies this package. Citations: the retained excerpts carry no citation command at all, so no bibliography entry of the article is transcribed and no reference is redistributed. Reference adaptation: none was needed and none was made; every reference inside the delivered unit resolves inside it, so no unresolved reference or citation is produced. Printed slips kept literal: no printed word, symbol, hypothesis or number of the counted statement or of its proof was altered. Layout and class: the article's own class is \texttt{elsarticle} with packages including amsmath, amsthm, amsfonts, amssymb, graphicx, colortbl, xcolor, float, subfig; the wrapper is the lane's pinned \texttt{amsart} editorial wrapper at 10pt on A4, which restates the theorem-like environments the retained text needs and adds the article's own macro definitions that the retained text needs (none), with extra wrapper packages (bm, bbm, dsfont, xspace, cancel, mathtools). The delivered numbering is the unit's own. Assets: no retained span contains a figure environment, a graphics-inclusion command, a PDF-page inclusion, a table or a TikZ picture; no image file is copied into this package, the package declares no asset dependency and no image file is redistributed with it. Licence: version arXiv:2505.16924v1 of this article is distributed under the Creative Commons Attribution 4.0 International licence, as recorded in the arXiv licence metadata for this exact version and shown on the abstract page https://arxiv.org/abs/2505.16924v1. A full rights notice is delivered at licenses/arxiv-2505.16924-CC-BY-4.0.txt. Characters: every retained excerpt is pure ASCII, so the delivered engine, XeLaTeX with its default Latin Modern text font, reports no missing character.
\begin{theorem}
\label{thm1}
For the operator $ T\in L^{A}(H) $ and $ q,\lambda,\mu\in \mathbb{C}, \ 0<\vert q \vert\leq 1, $ the following are true  \\
(1) $ \omega_{A,q}(T)\leq \Vert T \Vert_{A}. $ \\
(2) If  $ \alpha \in \mathbb{C}, \ \vert \alpha \vert =1, $ then $ \omega_{A,q}(\alpha T)=\omega_{A,\alpha q}(T). $ \\
(3) If  $ \alpha \in \mathbb{C}, \ \vert \alpha \vert =1,$ then $ c_{A,q}(\alpha T)=c_{A,\alpha q}(T). $ \\
(4) If $  \gamma =\sqrt{\vert \lambda\vert^{2}+ \vert \mu\vert^{2}+ 2Re (\lambda \overline{\mu} q)}\neq 0, $ then $ \gamma \omega_{A,\frac{\lambda + \mu \overline{q}}{\gamma}} (T)\leq \vert \lambda\vert\omega_{A}(T)+\vert \mu\vert\omega_{A,q}(T^{*}). $  \\
(5) If $  \gamma =\sqrt{\vert \lambda\vert^{2}+ \vert \mu\vert^{2}+ 2Re (\lambda \overline{\mu} q)}\neq 0, $ then $ \gamma c_{A,\frac{\lambda + \mu \overline{q}}{\gamma} }(T)\leq \vert \lambda\vert\omega_{A}(T)+\vert \mu\vert c_{A,q}(T^{*}). $\\
(6) For $ x,y \in H, \ \Vert x \Vert_{A}=\Vert y \Vert_{A}=1, \ \langle x,y \rangle_{A}=q, \ 0<\vert q \vert \leq 1 $, it is true $$
\Vert x\pm y \Vert_{A}=\sqrt{2}\sqrt{1\pm Req}.
$$ 
(7) $ \omega_{A}(T)\leq \omega_{A, q}(T)+ \sqrt{2(1-Req)}\omega_{A,\frac{1-q}{\sqrt{2(1-Req)}}}(T), \ q\neq 1. $ \\
(8) $ c_{A}(T)\leq c_{A, q}(T)+ \sqrt{2(1-Req)}\omega_{A,\frac{1-q}{\sqrt{2(1-Req)}}}(T), \ q\neq 1 . $
\end{theorem}

\begin{proof}
%\noindent {\it \textbf{Proof}} 
(1) For $ x, y \in H, $ $ \Vert x \Vert_{A}=\Vert y \Vert_{A}=1, \ \langle x,y \rangle_{A}=q, $ we have 
\begin{eqnarray}
\vert \langle Tx,y \rangle_{A}\vert & = & \vert \langle ATx,y \rangle \vert \nonumber \\
& \leq & \vert \langle ATx,Tx \rangle^{1/2} \vert \vert \langle Ay,y \rangle^{1/2} \vert \nonumber \\
& = & \vert \langle Tx,Tx \rangle_{A}^{1/2} \vert \vert \langle y,y \rangle_{A}^{1/2} \vert \nonumber \\
& = & \Vert Tx \Vert_{A} \Vert y \Vert_{A} \nonumber \\
& = & \Vert Tx \Vert_{A}.  \nonumber
\end{eqnarray}
Hence, it is obtained that
$$
\omega_{A,q}(T)\leq \Vert T \Vert_{A}.
$$
(2)-(3) These equalities are results of the following relations 
$$
\Vert \alpha x \Vert_{A}=\vert \alpha \vert \Vert  x \Vert_{A}=1,
$$
$$
\langle \alpha x,y \rangle_{A}=\alpha \langle x,y \rangle_{A}=\alpha q, \ 0<\vert \alpha q \vert= \vert q \vert \leq 1.
$$
(4)-(5) For the operator $ T\in L^{A}(H) $ and $ \lambda,\mu\in \mathbb{C},$ the following relation
\begin{eqnarray}
\vert \langle T(\lambda x+\mu y),x\rangle_{A} \vert &\leq & \vert \lambda\vert \vert \langle Tx,x\rangle_{A} \vert +  \vert \mu\vert \vert \langle Ty,x\rangle_{A} \vert  \nonumber \\
& = & \vert \lambda\vert \vert \langle Tx,x\rangle_{A} \vert +  \vert \mu\vert \vert \langle T^{*}x,y\rangle_{A} \vert \nonumber
\end{eqnarray}
is true. Since for $ \Vert x \Vert_{A}=\Vert y \Vert_{A}=1, \ \langle x,y\rangle_{A}=q, \ 0<\vert q \vert\leq 1, $ 
\begin{eqnarray}
\Vert \lambda x+\mu y \Vert_{A}^{2} & = & \langle A(\lambda x+\mu y), \lambda x+\mu y \rangle \nonumber \\
& = & \langle A(\lambda x), \lambda x \rangle +\langle A(\lambda x), \mu y \rangle +\langle A(\mu y), \lambda x\rangle + \langle A(\mu y), \mu y \rangle \nonumber \\
& = & \vert \lambda \vert^{2} \langle Ax,  x \rangle + \lambda \overline{\mu} \langle Ax,y \rangle + \overline{\lambda}\mu \langle Ay, x \rangle +\vert \mu \vert^{2} \langle Ay, y \rangle \nonumber \\
& = & \vert \lambda \vert^{2} \langle x,  x \rangle_{A} + \lambda \overline{\mu} \langle x,y \rangle_{A} + \overline{\lambda}\mu \langle y, x \rangle_{A} +\vert \mu \vert^{2} \langle y, y \rangle_{A} \nonumber \\
& = & \vert \lambda \vert^{2}+\vert \mu \vert^{2}+\lambda \overline{\mu}q+\overline{\lambda}\mu\overline{q}\nonumber \\
& = & \vert \lambda \vert^{2}+\vert \mu \vert^{2}+2Re\left( \lambda \overline{\mu}q \right), \nonumber
\end{eqnarray}
then for $ \gamma =\sqrt{\vert \lambda\vert^{2}+ \vert \mu\vert^{2}+ 2Re (\lambda \overline{\mu} q)} $ we have   $
\Vert \dfrac{\lambda x+\mu y}{\gamma}  \Vert_{A}=1 $   and $
\langle \frac{\lambda x+\mu y }{\gamma},x \rangle_{A}=\frac{\lambda +\overline{\mu} q }{\gamma}.$
Hence, from the above inequality and the definitions of $ A$-numerical and $ (A,q)$-numerical radii we have
$$
\gamma \omega_{A,\frac{\lambda + \mu \overline{q}}{\gamma}} (T)\leq \vert \lambda\vert\omega_{A}(T)+\vert \mu\vert\omega_{A,q}(T^{*}).
$$

Similarly, from the above inequality and the definitions of $ A$-numerical radius and $ (A,q)$-Crawford number it is implies
$$
\gamma c_{A,\frac{\lambda + \mu \overline{q}}{\gamma} }(T)\leq \vert \lambda\vert\omega_{A}(T)+\vert \mu\vert c_{A,q}(T^{*}).
$$
(6) For $ x,y \in H, \ \Vert x \Vert_{A}=\Vert y \Vert_{A}=1, \ \langle x,y \rangle_{A}=q, \ 0<\vert q \vert \leq 1 $, it is true 
\begin{eqnarray}\label{INeq_2}
\Vert x\pm y \Vert^{2}_{A} & = & \langle A(x\pm y),x\pm y \rangle \nonumber \\
& = & \langle Ax,x \rangle \pm \langle Ax,y \rangle \pm \langle Ay,x \rangle + \langle Ay,y \rangle \nonumber \\
& = & \Vert x \Vert_{A}^{2}\pm 2 Re \langle x,y \rangle_{A}+ \Vert y \Vert_{A}^{2} \nonumber \\
& = & 2 \pm 2 Req. 
\end{eqnarray}
Hence, we have
$$
\Vert x\pm y \Vert_{A}=\sqrt{2}\sqrt{1\pm Req}.
$$
(7) For any $ x,y \in H, \ \Vert x \Vert_{A}=\Vert y \Vert_{A}=1, \  \langle x,y \rangle_{A}=q, \ q\neq 1, $ we have 
\begin{eqnarray}
\vert \langle Tx,x\rangle_{A} \vert & \leq & \vert \langle Tx,y\rangle_{A} \vert +  \vert \langle Tx,x-y\rangle_{A} \vert \nonumber \\
& = & \vert \langle Tx,y\rangle_{A} \vert + \sqrt{2(1-Req)}\vert \langle Tx,\frac{x-y}{\sqrt{2(1-Req)}}\rangle_{A} \vert . \nonumber 
\end{eqnarray}
Then, from (6), we have $ \Vert \frac{x-y}{\sqrt{2(1-Req)}} \Vert_{A}=1 $ and $ \langle x, \frac{x-y}{\sqrt{2(1-Req)}} \rangle =\frac{1-q}{\sqrt{2(1-Req)}}. $
Hence, from the previous relation and the definitions of $ A$-numerical and $ (A,q)$-numerical radii we have
$$
\omega_{A}(T)\leq \omega_{A, q}(T)+ \sqrt{2(1-Req)}\omega_{A,\frac{1-q}{\sqrt{2(1-Req)}}}(T).
$$
(8) This claim can be proved similarly. \\ \\
  \end{proof}

\end{document}
